Prostorová distribuce \aclgen{TFR} v~\acs{geo-EU27} v~roce~2024 vykazuje rozsah od~\num{1,18} (Malta) po~\num{1,79} (Bulharsko); medián činí~\num{1,42}. Nejvyšší hodnoty v~severozápadním klastru: Francie~\num{1,75}, Švédsko~\num{1,52}, Dánsko~\num{1,55}, Irsko~\num{1,57}.

Středovýchodní pás (\ac{ČR}~\num{1,45}, Slovensko~\num{1,49}, Rakousko~\num{1,44}, Slovinsko~\num{1,52}) leží v~pásmu \numrange{1,40}{1,55}. Jižní okraj (Malta~\num{1,18}, Itálie~\num{1,21}, Španělsko~\num{1,16}) tvoří klastr nejnižší plodnosti. Hodnota \aclgen{ČR} odpovídá průměru \acs{geo-EU27} (\num{1,46}).

\inputpgffigure{vyhled_porodnost_mapa}
